\subsection{Bishop operators: exact finite certificates for the cyclicity question}
\label{dev:bishop}

\paragraph{The question supported by the notebook.}
On P000089 the notes give, on complex \(L^2(0,1)\),
\[
 (B_\alpha f)(x)=x f(\{x+\alpha\}),\qquad
 \alpha\in(0,1)\setminus\mathbb Q,
\]
where \(\{\cdot\}\) denotes fractional part.  The margin says
\ru{Изучить подробно этот оператор}.  The text then juxtaposes the suggestion
\ru{нет инв. подпр.}\ with the admission
\ru{неизв. ни одного циклич. вектора}.  These are different quantifiers.
For a bounded operator \(B\), absence of a nonzero proper closed invariant
subspace is equivalent to \emph{every nonzero} vector being cyclic; existence
of \emph{one} cyclic vector is much weaker.  Indeed, the closed linear span
of the orbit of a nonzero vector is an invariant subspace, and conversely a
vector inside a proper invariant subspace cannot be cyclic.  The zero vector
must be excluded from the notebook's shorthand \(\forall\) assertion.

The precise questions are therefore: for which irrational \(\alpha\) does
\(B_\alpha\) have a nonzero proper closed invariant subspace, and for which
does it possess a cyclic vector?  The stronger, plausible cyclicity question
considered here asks whether the \emph{specific} vector \(1\) is cyclic for
every irrational \(\alpha\).  This last universal formulation is an editorial
research question, not a statement attributed to Halmos or Gelfand.
No universal answer is claimed below.

\paragraph{Verified literature and the scope of the development.}
\citet[Theorem~3.7]{chamizo2020bishop}, first circulated on 19 December 2018
and published in 2020, prove the existence of nontrivial hyperinvariant
subspaces if the continued-fraction denominators satisfy
\[
 \log q_{n+1}=O\!\left(\frac{q_n}{(\log q_n)^3}\right).
\]
Their Theorem~4.1 supplies an obstruction to the particular Beurling-algebra
method for sufficiently extreme approximability; it does not prove absence
of invariant subspaces.  \citet[Theorems~4.4, 4.5, 4.9, 4.16]{behani2023bishop},
first posted on 23 July 2022 and published in 2023, give the rational
fibre-determinant criterion, rational cyclicity of holomorphic vectors
nonzero at zero, generic irrational cyclicity, and a sufficient large-gap
condition.  Locators here refer to the opened arXiv version.

A search dated 12 November 2026 also found the abstract of
\citet{ma2026rational}, which announces quantitative determinant and cyclicity
results and specifically mentions \((r,q)=(9,16)\).  Its primary full-text
and abstract endpoints did not open in this run.  It is marked
\emph{unverified}; none of its announced theorems is used as an assumption.
Its existence prevents a novelty claim for the same finite example or for
the general programme of explicit gap estimates.  No all-irrational
cyclicity or invariant-subspace theorem was verified in the sources checked
through 12 November 2026.

The bounded contribution is an exact-arithmetic certificate format, an
independent proof of the transfer estimate, and a fully specified terminating
procedure for a sufficient gap function.  The rational and generic
cyclicity conclusions are established literature, re-proved here for the
certificate's foundations.  The label \ours{} means that this companion
supplies and checks a proof; it does not assert historical priority.

\paragraph{Rational fibres and the endpoint determinant.}
Fix relatively prime integers \(1\le r<q\), \(q\ge2\), and write
\(a_i=(ir\bmod q)/q\), with indices read modulo \(q\).
Splitting \([0,1)\) into the intervals \(a_i+[0,1/q)\) identifies
\(L^2(0,1)\) unitarily with
\(L^2((0,1/q);\mathbb C^q)\) by
\[
 (Uf)(t)_i=f(t+a_i),\qquad 0<t<1/q.
\]
Lebesgue measure is used on the base and the ordinary Euclidean norm on the
fibres, so no normalization factor is missing.  In these coordinates,
\[
 (\mathcal B(t)v)_i=(t+a_i)v_{i+1},\qquad
 \mathcal B(t)^q=m_q(t)I,
 \quad m_q(t)=\prod_{j=0}^{q-1}\left(t+\frac jq\right).
\]
The continuous function \(m_q\) is strictly increasing on \([0,1/q]\).
Let \(K_{r,q}(t)\) be the matrix whose column \(k\), \(0\le k<q\), is
\(\mathcal B(t)^k\mathbf1\), and let \(D_{r,q}(t)=\det K_{r,q}(t)\).

\begin{lemma}[A nonzero endpoint determinant]\ours{}
\label{lem:bishop-endpoint}
For the coprime pair above,
\[
 K_{r,q}(t)_{i,k}=\prod_{h=0}^{k-1}(t+a_{i+h}),\qquad
 D_{r,q}(0)=(-1)^{(q-1)(q-2)/2}\prod_{j=1}^{q-1}a_j^{\,j}\ne0.
\]
An empty product is one.  Consequently \(D_{r,q}\) is a nonzero polynomial
and has only finitely many zeros in the fundamental interval.
\end{lemma}
\begin{proof}
The product formula follows by iterating the displayed weighted cyclic
matrix.  Since \(a_0=0\) and every other \(a_i\) is positive, row zero at
\(t=0\) is \((1,0,\ldots,0)\).  In row \(i\ge1\), the last possibly
nonzero entry is column \(q-i\), equal to
\(\prod_{j=i}^{q-1}a_j\).  Delete the first row and column and reverse the
remaining \(q-1\) rows.  The resulting triangular matrix has these entries
on its diagonal.  Reversal has sign
\((-1)^{(q-1)(q-2)/2}\); multiplying the diagonal entries counts \(a_j\)
exactly \(j\) times.  This proves the formula and all its consequences.
\end{proof}

\begin{proposition}[Rational cyclicity of the constant vector]\ours{}
\label{prop:bishop-rational}
The vector \(1\) is cyclic for \(B_{r/q}\) on complex \(L^2(0,1)\).
\end{proposition}
\begin{proof}
Let \(\mathcal C\) be the closed span of \(B_{r/q}^n1\), \(n\ge0\),
in the fibre representation.  For each \(0\le k<q\) it contains
\(p(m_q(t))\mathcal B(t)^k\mathbf1\) for every polynomial \(p\), since
\(B_{r/q}^{q\ell+k}1\) has precisely that form for \(p(z)=z^\ell\).
The map \(m_q\) is a homeomorphism onto a compact real interval.
Polynomial approximation therefore gives every continuous scalar
coefficient \(h(t)\) times this bounded vector field.  Approximation of
bounded measurable coefficients by continuous functions in the appropriate
finite weighted \(L^2\) measure gives the same conclusion for bounded
measurable \(h\).

For arbitrary \(g\in L^2((0,1/q);\mathbb C^q)\), define
\[
 E_N=\{t:\|K_{r,q}(t)^{-1}\|\le N,\ \|g(t)\|\le N\},
 \qquad g_N=\mathbf1_{E_N}g.
\]
The finitely many singular fibres can be discarded.  By the lemma,
\(E_N\) exhausts the base up to a null set; hence \(g_N\to g\) in
\(L^2\).  The components of
\(\mathbf1_{E_N}K_{r,q}^{-1}g\) are bounded measurable functions.
Expressing \(g_N\) in the columns of \(K_{r,q}\) puts \(g_N\) in
\(\mathcal C\).  Closedness now gives \(g\in\mathcal C\).
\end{proof}

\paragraph{Why a uniformly invertible fibre matrix is the wrong requirement.}
The proof needs almost-everywhere invertibility and truncation, not a
positive lower bound for \(|D_{r,q}(t)|\) on the whole interval.  This
distinction is testable without estimating any irrational orbit.

\begin{example}[A singular interior fibre]\oursComputed{}
\label{ex:bishop-singular}
For \((r,q)=(9,16)\), exact arithmetic gives
\(D_{9,16}(0)<0<D_{9,16}(1/16)\).  Thus at least one
\(t\in(0,1/16)\) has \(D_{9,16}(t)=0\), although the constant vector
is cyclic by Proposition~\ref{prop:bishop-rational}.
\end{example}
For reproducibility, use the integer matrix
\[
 C(y)_{i,k}=\prod_{h=0}^{k-1}\bigl(y+(9(i+h)\bmod16)\bigr),
 \qquad \det C(y)=16^{120}D_{9,16}(y/16).
\]
The two exact determinants, including their integer signs, are in
\nolinkurl{evidence/independent_bishop_check.json}.  The first computation
used rational symbolic Gaussian elimination; the second used a separately
implemented fraction-free integer Bareiss algorithm, with every division
checked for exactness.  Both agree.  Continuity and the intermediate value
theorem, rather than a plotted sample, establish the existence of the zero.
No assertion is made about its uniqueness, multiplicity, or the location
of zeros for all other coprime pairs.

For small denominators the same program obtains, for example,
\[
 D_{1,2}(t)=\frac12,\quad
 D_{1,3}(t)=-\frac{9t+4}{27},\quad
 D_{2,3}(t)=-\frac{9t+2}{27}.
\]
Checks of the endpoint formula and matrix-power identity for all 31 coprime
pairs with \(2\le q\le10\) passed.  Those checks are finite checks;
Lemma~\ref{lem:bishop-endpoint} is the independent all-\(q\) proof.

\paragraph{A quantitative transfer bound that accounts for wrap points.}
For \(k\ge0\), put \(w_{\alpha,k}=B_\alpha^k1\).  Thus
\(w_{\alpha,0}=1\) and
\[
 w_{\alpha,k}(x)=\prod_{j=0}^{k-1}\{x+j\alpha\}\quad(k\ge1).
\]
All factors lie in \([0,1]\).  Treating them as globally Lipschitz
functions of \(\alpha\) would overlook their jumps at wrap points.

\begin{lemma}[Finite-orbit continuity]\ours{}
\label{lem:bishop-continuity}
For all real \(\alpha,\beta\) and integers \(k\ge0\),
\[
 \|w_{\alpha,k}-w_{\beta,k}\|_2
 \le |\alpha-\beta|^{1/2}\sum_{j=1}^{k-1}\sqrt j
 \le \frac{k(k-1)}2|\alpha-\beta|^{1/2},
\]
where the sum and final bound are zero for \(k=0,1\).
Consequently, if \(Q(z)=\sum_{k=0}^N c_kz^k\), then
\[
 \|Q(B_\alpha)1-Q(B_\beta)1\|_2
 \le R(Q)|\alpha-\beta|^{1/2},\qquad
 R(Q)=\sum_{k=0}^N|c_k|\frac{k(k-1)}2.
\]
\end{lemma}
\begin{proof}
For \(0\le d\le1\), a sawtooth translated by \(d\) changes by
\(d\) on an interval of length \(1-d\) and by \(d-1\) on the remaining
interval.  Its squared \(L^2\) difference is
\((1-d)d^2+d(1-d)^2=d(1-d)\).  The same identity applies to a
negative shift by replacing it by its absolute value modulo one.
Translation invariance then gives
\[
 \|\{x+j\alpha\}-\{x+j\beta\}\|_2^2
 \le j|\alpha-\beta|.
\]
For \(j|\alpha-\beta|\le1\), use the identity; for a larger difference,
the identity is at most \(1/4\le j|\alpha-\beta|\).
Expand the difference of the two products by telescoping.  Every other
factor has absolute value at most one, so the triangle inequality in
\(L^2\) gives the first bound.  For integers \(j\ge1\),
\(\sqrt j\le j\), giving the second.  Another triangle inequality,
now over the coefficients of \(Q\), proves the final claim.
\end{proof}

\begin{example}[A complete finite approximation certificate]\oursComputed{}
\label{ex:bishop-certificate}
Let
\[
\begin{split}
 Q(z)=\frac{1}{1032182056}\bigl(&-24558043+609112200z
       -604809944z^2\\
       &+2030693280z^3-987436800z^4\bigr).
\end{split}
\]
Direct piecewise polynomial integration gives
\[
 \|Q(B_{1/2})1-x^2\|_2^2
 =\frac{2117561}{24772369344}<\frac1{10000},
 \qquad R(Q)=\frac{1577688823}{129022757}<13.
\]
Therefore, for every real \(\alpha\),
\[
 |\alpha-\tfrac12|<\frac1{1690000}
 \quad\Longrightarrow\quad \|Q(B_\alpha)1-x^2\|_2<\frac1{50}.
\]
This is a statement about this target and this accuracy, valid for rational
and irrational parameters in the indicated neighbourhood.
\end{example}
Indeed, Lemma~\ref{lem:bishop-continuity} bounds the additional error by
\(13/1300=1/100\), and the rational approximation error is strictly less
than \(1/100\).  It does not follow from a single such certificate that
\(1\) is cyclic for any fixed irrational parameter.

The coefficients were obtained by exact least squares in the span of the
first five orbit vectors.  One program used the normal equations and then
integrated the squared residual independently.  A second program used only
Python's rational-number arithmetic, rebuilt the orbit polynomials, and
reintegrated both cells.  It obtained the identical rational number.
The same two checks agree on all five certificates in the output:
\begin{center}
\begin{tabular}{@{}rrrcc@{}}
\toprule
\(r\)&\(q\)&Degree&\(L^2\) error for \(x^2\)&\(R(Q)\)\\
\midrule
1&2&2&0.07448495&0.02173913\\
1&2&4&0.009245581&12.22799\\
1&2&8&0.0007486373&332.0757\\
1&3&6&0.01199983&91.78700\\
2&3&6&0.008914569&58.38182\\
\bottomrule
\end{tabular}
\end{center}
These rounded values describe exact fractions saved with the code.  The
tradeoff is visible: a smaller rational error can require a much larger
coefficient sum, reducing the neighbourhood to which the certificate
transfers.  Small residuals alone do not control irrational parameters.

\paragraph{A terminating sufficient criterion, with all quantifiers stated.}
The finite certificates can be organized into a theorem.  This produces an
algorithmically computable function; no modest closed-form growth bound for
it is asserted.

\begin{theorem}[A computable continued-fraction gap rule]\ours{}
\label{thm:bishop-gap}
There is a total computable function \(\Psi:\{2,3,\ldots\}\to\mathbb N\)
with the following property.  If an irrational \(\alpha\in(0,1)\), with
continued-fraction convergents \(r_n/q_n\), satisfies
\(q_{n+1}>\Psi(q_n)\) for infinitely many \(n\) with \(q_n\ge2\), then
the vector \(1\) is cyclic for \(B_\alpha\) on complex \(L^2(0,1)\).
\end{theorem}
\begin{proof}
For each fixed \(q\), each \(1\le r<q\) coprime to \(q\), and each
\(1\le j\le q\), enumerate polynomials \(Q\in\mathbb Q[z]\) in a fixed
effective order until
\[
 \|Q(B_{r/q})1-x^{j-1}\|_2^2<\frac1{4q^2}.
 \tag{*}
\]
This is a decidable strict rational inequality: split \([0,1)\) at the
points \(i/q\), express the orbit products as rational polynomials on
each cell, and integrate their squared residual.  To specify the enumeration,
at stage \(h=1,2,\ldots\) list degrees \(0,\ldots,h\) and all coefficient
tuples \(c_k=a_k/b_k\), \(|a_k|\le h\), \(1\le b_k\le h\), in
lexicographic order of the integer tuples; repetitions are harmless.
Every rational polynomial eventually occurs.  The search terminates
by Proposition~\ref{prop:bishop-rational}.  Since the target and the orbit
vectors are real, take real parts of an approximating complex polynomial
and then approximate its finitely many coefficients by rationals; the
strict target error remains attainable.

Call the first polynomial so found \(Q_{r,q,j}\), and define the positive
rational number
\[
 A_q=\max_{\substack{1\le r<q,\ (r,q)=1\\1\le j\le q}}
       \bigl(1+R(Q_{r,q,j})\bigr),
 \qquad \Psi(q)=\lceil4q A_q^2\rceil.
\]
Only finitely many terminating searches occur for each input \(q\), so
\(\Psi\) is total and computable.  The usual convergent inequality is
\(|\alpha-r_n/q_n|<1/(q_nq_{n+1})\).  On the stipulated subsequence it
therefore gives
\[
 |\alpha-r_n/q_n|<\frac1{4q_n^2 A_{q_n}^2}.
\]
Combining (*) with the continuity lemma yields, for every \(j\le q_n\),
\[
 \|Q_{r_n,q_n,j}(B_\alpha)1-x^{j-1}\|_2<\frac1{q_n}.
\]
The denominators on any infinite subsequence tend to infinity.  For each
fixed \(j\), the monomial \(x^{j-1}\) thus belongs to the closed cyclic
subspace.  Polynomials are dense in \(L^2(0,1)\), proving cyclicity.
\end{proof}

\begin{corollary}[The fixed-vector generic conclusion]\ours{}
The set of irrational \(\alpha\in(0,1)\) for which \(1\) is cyclic for
\(B_\alpha\) contains a dense \(G_\delta\) set.
\end{corollary}
\begin{proof}
For each positive pair \(j,m\), take the union, over rational-coefficient
polynomials \(Q\), of the open sets
\(\{\alpha:\|Q(B_\alpha)1-x^{j-1}\|_2<1/m\}\).
Openness follows from Lemma~\ref{lem:bishop-continuity}; each set in the
countable family of unions contains every rational \(\alpha\in(0,1)\)
by Proposition~\ref{prop:bishop-rational}, so is dense.  The Baire theorem
and intersection with the irrational parameters give the conclusion.
All monomials are approximable at the parameters in this intersection.
\end{proof}

\paragraph{Independent proof check and the remaining obstruction.}
The fibre proof was re-derived using the direct cyclic matrix rather than
the orbit-product recurrence; the factors and row-reversal sign agree.
The continuity bound was checked separately by splitting at the sawtooth
jump and obtaining \(d(1-d)\); a pointwise Lipschitz bound was rejected
because it fails at that jump.  The gap proof was re-checked in the order
fixed target, infinite convergent subsequence, denominator tending to
infinity, and finally polynomial density.  The small cases \(k=0,1\) and
\(q=2\) are covered explicitly.  The two independent exact programs agree
on the determinant signs and every reported integral, including both
numerator orders at denominator three.  No finite-denominator observation
has been promoted into an all-denominator assertion.

The certificate method lacks a quantitative bound on the degree and
coefficient size needed as \(q\) grows.  Even a practical bound would only
give a sufficient approximation class unless the final transfer argument
were strengthened.  It also cannot turn existence of one cyclic vector
into cyclicity of every nonzero vector or settle the invariant-subspace
question.  Determining the full scope of arbitrary-weight classification
criteria and addressing the unrestricted Hilbert invariant-subspace
problem lie outside this bounded development; the ambiguous walk question
does not yet specify a unique probability law.

One further research run should first obtain and check the 2026 determinant
preprint, compare its quantitative bounds with these exact certificates,
and then optimize the rational approximation step with proved coefficient
bounds.  It could supply a certified comparison on additional finite
denominators and possibly a sharper sufficient arithmetic condition.
Neither those computations nor this chapter justify a forecast of a
universal irrational theorem.
